\chapter{Introduzione}
\section{Metodi cose}

\[
  \|u - u_x\|_V \le \frac{M}{\alpha} \mathrm{inf}_{v_k \in V_k} \| u - v_k\|_V
\]

\begin{example}{}
    Consideriamo il problema di Poisson 
    \[
      \begin{cases}{}
          - \mathrm{div}{(k \nabla u)} = f\\
          u|_{\mathcal{U}} = 0
      \end{cases}
    \]
    dato da 
    \[
      a {(w, v )} = \int _{\Omega} k{(x_{1}, x_{2})} \nabla w {(x_{1}, x_{2})}
      \cdot \nabla v {(x_{1}, x_{2})} \,dx_{1}\, dx_{2}
    \]
    con \(V = H^{1}_0{(\Omega)}\).

    Allora abbiamo che \(a\) è continua:
    \[
      | a{(w, v)} | = \left| \int_\Omega \dots \right| \le
      \|k\|_{L^{\infty}{(\Omega)}}  \, |w|_{H^{1}{(W)}} \, | v | H^{1}{(W)}
    \]
    mentre per la coercività:
    \[
      a{(w, w)} = \int _\Omega k | \nabla w | ^2 \ge k_\text{min} |w|^2_{H^{1}}
      \ge \underbrace{\frac{k_\text{min} }{2} \mathrm{min}\left\{1,
      \frac{1}{c^2_p}\right\}}_{\alpha} 
      \|w\|^2_{H^{1}} 
    \]
    Dunque abbiamo che
    \[
      \frac{M}{\alpha} = C \frac{k_\text{max} }{k_\text{min} } \gg 1
    \]
\end{example}

\begin{example}{}
Sia \(u\) la temperatura e \(\mathbf{q} \) il vettore che indica il flusso di
calore. Abbiamo che il calore ``va dal caldo al freddo'', dunque
\[
  \mathbf{q}  = - \kappa \nabla u
\]
con \(\kappa\) la \emph{conducibilità} del materiale che è una funzione della
posizione. Notare che \(\kappa\) potrebbe anche essere una matrice, se il
materiale è \emph{anisotropo}, perché il calore non conduce in tutte le
direzioni ugualmente. Un esempio sono i materiali compositi di fibre e colla.

A questo punto in ogni pezzettino di pentola \(\omega \subseteq \Omega \)
abbiamo che:
\[
  \int_{\partial \omega} \mathbf{q} \cdot \mathbf{n}  - \int_{\omega} f = 0
\]
per il teorema di Gauss abbiamo che \(\int_\omega f = - \int _{\omega}
\mathrm{div} k \nabla u\). Se supponiamo che l'uguaglianza valga per ogni
\(\omega\), allora abbiamo che su tutto \(\Omega\) 
\[
  \begin{cases}{}
      - \mathrm{div} k \nabla u = f & \in \Omega \\    
      u|_{\Gamma_{0}} & \text{è assegnata} \\
      - (k \nabla u_{\Gamma_n} \cdot \mathbf{n} )& \text{è assegnato}
  \end{cases}
\]
l'ultima, che è la condizione di Neumann, ad esempio se fosse assegnato a 0
significherebbe che non conduce calore in alcuni punti?
\end{example}

Se \(a{(\cdot , \cdot )}\) è anche simmetrica, allora \(\|w\|^2_E = a{(w, w)}
\ge \alpha \|w\|^2_V\) e notando che \(a{(w, w)} \le M \|w\|^2_V\) allora
notiamo che \(\|\cdot \|_E\) e \(\|\cdot \|_V\) sono equivalenti.

Allora per il metodo di Galerkin abbiamo che \(a{(w, w)} \ge \|w\|^2_E\) e
\(a{(w, v)} \le \|w\|_E \|v\|_E\). Segue che
\[
  \sqrt{\alpha} \|u - u_k\| \le  \|u - u_k\|_E \le \inf_{v_n \in V_n}
  \|u - v_n\|_E \le \sqrt{M} \inf \|u - v_n\|_V 
\]
da cui infine 
\[
  \|u - u_n\|_E \le \frac{\sqrt{M}}{\sqrt{\alpha}} \inf_{v_{n} \in V_n} \|u -
  v_n\|_V
\]
che migliora l'errore.

\section{Metodo degli Elementi Finiti (FEM)}

Per trovare \(V_k\) nel caso dei FEM. 

\begin{enumerate}[label = \arabic*.]
    \item \(\Omega \approx \Omega_k = \cup_{T \in  \mathcal{T}_h} T\) 
    \item \(V_k\) sono poligoni 
\end{enumerate}

La triangolazione è \(\mathcal{T}_h = \{T_{1}, \dots, T_{N_\text{el}} \} \) ed è
tale che \(\cup_{T \in \mathcal{T}_n} T = \overline{\Omega}\). Perché
\(\mathcal{T}_k\) sia ammissibile imponiamo che per ogni \(i, j\), \(T_{i}
\cap T_j\) può solo essere un triangolo intero, un lato intero, un vertice o
nulla.

Preso ora \(T \in \mathcal{T}_h\), definiamo due valori \(h_T = \mathrm{diam}\,T\)
e \(\rho_T\) il raggio del più grande cerchio inscrivibile. Allora diciamo che
la famiglia di triangolazioni \(\{\mathcal{T}_h\}_{h > 0}  \) se esiste una
\(C_s \in \mathbb{R}\) tale che \(\forall h > 0\) e \(T \in \mathcal{T}_h\),
abbiamo che \(\frac{h_T}{\rho_T} \le C_s\). Chiaramente le stime di errore
avranno all'interno la costante \(C_s\). Una buona famiglia di triangolazioni è
infatti tale che i triangoli tendano a essere circa punti, e non rette o
segmenti.

